%----------------------------------------------------------------------------
\chapter{Related Work and Gap Analysis}\label{ch:related}
%----------------------------------------------------------------------------
\guidetodo{Related Work and Gap Analysis}{app:chapters:related}

\section{Comparison criteria}
\guidetodo{Comparison criteria}{app:chapters:comparison-criteria}

\section{Existing approaches}
\guidetodo{Existing approaches}{app:chapters:existing-approaches}
TODO The related work by family; \cref{fig:positioning} places it, and \cref{tab:related} compares the closest approaches.

\begin{runningexample}[elsewhere]
TODO How the closest approach would handle the scenario, and which requirement it fails on.
\end{runningexample}

\begin{figure}[!ht]
  \centering
  \missingfigure[figheight=2cm]{Positioning figure: \cref{fig:guide:positioning}}
  \caption{TODO positioning figure.}
  \label{fig:positioning}
\end{figure}

\begin{table}[!ht]
  \centering
  \caption{TODO Gap matrix: each requirement against the closest approaches. \gaplegend.}
  \label{tab:related}
  \begin{tabularx}{\linewidth}{l X c c c c}
    \toprule
    & \textbf{Requirement} & \makecell{\textbf{Approach A}\\\cite{lamport1978time}} & \makecell{\textbf{Approach B}\\\cite{w3c2025vcdm}} & \textbf{Approach C} & \textbf{This work} \\
    \midrule
    \ref{req:first} & TODO short name & \gapyes & \gappartial & \gapno & \gapyes\,(\cref{rq:correct}) \\
    \bottomrule
  \end{tabularx}
\end{table}

\section{Gap}\label{sec:gap}
\guidetodo{Gap}{app:chapters:gap}
